\section{Discussion}
\label{sec:limitations}

BDL is a training objective for sequence models evaluated by ranked temporal detections.
It changes what the model is asked to learn, but it does not replace the rest of the detector.
A trained model still produces score sequences, and an application still chooses peak selection, thresholds, and any validity constraints.
Its purpose is to align training with the temporal decisions later matched by AP.
Segmentation asks the model to label samples, while BDL asks it to estimate the expected number of events in each output bin.
That distinction covers point events, onset-only labels, repeated events,
interval endpoints, and multi-type streams without requiring the likelihood to
encode interval structure.
If an application needs valid intervals, duration priors, or endpoint pairing,
those constraints remain decoder choices.
If deployment consumes calibrated labels at every sample, segmentation remains
the direct objective.

The sleep benchmark exposes the training--evaluation mismatch especially clearly because long annotated intervals are scored by much shorter tolerance windows around onset and wake-up times.
Across a recurrent encoder and U-Net variants, the improvement is largest where event AP depends most on the ranking of localized event scores.
Class-weighted and focal segmentation losses address imbalance, but they still train on interval-mask labels and recover event scores afterward.
Their lower performance suggests that the issue is not only class prevalence, but also the mismatch between state-occupancy training and event-time evaluation.

BDL also separates event-localized training from generic localized labels.
Smoothing can help when annotation times are uncertain, but an annotation
should still represent one event after clipping, downsampling, or overlap
with nearby events.
In BDL, smoothing redistributes that event over nearby times and binning sums
the contributions onto the model timeline.
The output is therefore the expected number of events in each bin, giving the
sequence a fixed unit where proximity maps or normalized location distributions
would leave scale as a modeling convention.
The state-reconstruction analysis makes the complementary point: learned
score sequences can be integrated back into state sequences, including through
a hazard recursion that treats onset and wake-up scores as transition
probabilities over time.
Direct segmentation remains stronger when performance is measured by
samplewise state-label accuracy rather than ranked detections.

The seizure and streaming-context ablations expose the remaining dependencies.
The smoothing kernel encodes annotation uncertainty, the receptive field must
include the evidence that makes an event recognizable, and some domains need
calibrated thresholds or structured decoders.
BDL specifies what the sequence output should estimate, but architecture,
context, calibration, and post-processing determine whether the learned scores
become accurate detections for a particular signal.
This is the intended scope of the method: it supplies a statistically defined
training signal for event detection while leaving representation learning and
domain constraints as detector design choices.
